\chapter{Introduction and Problem Formulation}
\label{ch:introduction}

\section{Institutional Context and the Examination Grading Bottleneck}
\label{sec:institutional_context}

At premier engineering institutes like the Indian Institute of Technology Guwahati (IIT Guwahati), undergraduate mathematics education still centers around handwritten pen-and-paper examinations. In large foundational courses such as \textbf{MA102: Mathematics-II (Linear Algebra and Multivariable Calculus)}, offered by the Department of Mathematics, enrollment regularly exceeds 900 first-year undergraduate students in a single semester.

Under standard academic procedures, each student writes their solutions in an unruled 14-page physical answer booklet. Over an entire semester---including two quizzes, the mid-semester exam, and the end-semester exam---this creates an enormous volume of scanned paper: over $\mathbf{13,146}$ pages per exam cycle. Grading this mountain of handwritten mathematics creates severe practical difficulties for instructors and teaching assistants:
\begin{enumerate}
    \item \textbf{Cognitive Fatigue and Inconsistent Grading:} Grading hundreds of pages of unconstrained, hurried handwriting over several days is physically exhausting. Evaluators get tired, and their grading standards can easily drift between the first booklet they read in the morning and the last one late at night.
    \item \textbf{Dense Two-Dimensional Notation:} Unlike humanities or essay exams, mathematics scripts are packed with two-dimensional structures. A single page might contain a $4 \times 4$ real matrix undergoing elementary row operations ($R_3 \to R_3 - \frac{1}{2} R_1$), followed by a multivariable calculus integral ($\iint_D f(x, y) \, dx dy$) or an inductive proof of a vector subspace ($c\mathbf{u} + \mathbf{v} \in W$).
    \item \textbf{Messy Unruled Writing and Crossed-Out Work:} Because students write on blank unruled paper under time pressure, their writing drifts diagonally across the page. Lines overlap, pen pressure varies, ink bleeds through the paper, and students frequently scratch out mistakes with heavy scribbles before restarting their derivations.
\end{enumerate}

If we can automatically convert these raw scanned exam pages into neat, compilable \LaTeX\ source code, evaluators could review clean typeset mathematics on their screens alongside the original scanned page. This would eliminate the headache of deciphering messy handwriting while keeping the human instructor firmly in charge of assigning marks.

\section{Mathematical Typesetting and First Principles}
\label{sec:first_principles_latex}

Before designing an automated transcription pipeline, we first need to look closely at what \LaTeX\ actually is and how it represents mathematical thought.

Created by Leslie Lamport on top of Donald Knuth's \TeX\ engine \cite{knuth1984texbook}, \LaTeX\ is not just a digital document format or a simple stream of text characters. It is a formal compiled macro language. Knuth modeled typesetting using a physical metaphor of \emph{boxes} (rigid rectangular blocks with specific widths, heights, and depths) and \emph{glue} (flexible spacing that stretches or shrinks to balance lines).

In plain text, each character follows the previous one in a simple line. In contrast, \LaTeX\ treats a document as an \textbf{Abstract Syntax Tree (AST)}. Consider a simple fraction like $\frac{\alpha + \beta}{\det(\mathbf{A})} = \lambda_1$. In \LaTeX, the command \texttt{\textbackslash frac} tells the compiler to construct two separate branches in the syntax tree: one box for the numerator above the horizontal line, and one box for the denominator below it. 

\subsection{Why Mathematicians Rely on \LaTeX}
\LaTeX\ has remained the gold standard in mathematics and physics for four decades because it enforces structural precision:
\begin{enumerate}
    \item \textbf{Clear Semantic Roles:} It separates scalars ($x$), bold vectors ($\mathbf{x}$), limits ($\lim_{x \to 0}$), exponents ($x^2$), and index subscripts ($x_2$).
    \item \textbf{Two-Dimensional Grids:} Linear algebra relies on matrix environments (\texttt{\textbackslash begin\{bmatrix\} \dots \textbackslash end\{bmatrix\}}) that align rows and columns in a rigorous 2D grid.
    \item \textbf{Zero Syntax Tolerance:} A \LaTeX\ source file must be mathematically and syntactically watertight. If an author forgets even a single closing brace (\texttt{\}}) or misspells an environment delimiter, the \TeX\ compiler aborts immediately with a fatal syntax error.
\end{enumerate}

\subsection{The Core Dimensionality Mismatch: 1D Text vs. 2D Math}
When we started our research, our first key realization was that reading printed text and reading handwritten math are fundamentally different problems. We summarize this contrast in Table \ref{tab:1d_vs_2d_ocr}.

\begin{table}[htbp]
\centering
\caption{First-Principles Comparison: 1D Printed Text OCR vs. 2D Handwritten Math HMER}
\label{tab:1d_vs_2d_ocr}
\vspace{0.2cm}
\small
\begin{tabular}{p{2.8cm}p{5.6cm}p{5.6cm}}
\toprule
\textbf{Structural Dimension} & \textbf{Printed Text OCR (1D Sequence)} & \textbf{Handwritten Math HMER (2D Graph)} \\
\midrule
\textbf{Spatial Topology} & Left-to-right linear reading sequence. & Non-linear 2D planar layout; nested tree hierarchy. \\
\addlinespace
\textbf{Baseline Geometry} & Straight, horizontal printed baselines. & Curving, drifting baselines; variable pen slant. \\
\addlinespace
\textbf{Symbol Meaning} & Mostly independent of vertical position. & Strictly depends on vertical offset ($x^2$ vs. $x_2$). \\
\addlinespace
\textbf{Scope Delimiters} & Simple commas, periods, quotation marks. & Multi-line brackets, square roots, fraction bars. \\
\addlinespace
\textbf{Current State} & Highly reliable ($>99\%$ character accuracy). & Open research problem ($\approx 50$--$65\%$ on whole exam pages). \\
\bottomrule
\end{tabular}
\end{table}

Standard OCR systems scan text in a linear left-to-right pass. But math notation is a two-dimensional graph that requires extracting a \textbf{Spatial Relationship Tree (SRT)}. In an exam script, a fraction bar splits the space into numerator and denominator sub-regions; exponents float upward and shrink in size; and matrices form 2D grids where entries align both horizontally and vertically. When classical 1D text algorithms encounter this 2D spatial layout, their internal line-by-line assumptions break down completely.

\section{Why Existing Systems Struggle with Handwritten Math Exams}
\label{sec:failure_modes}

Although optical character recognition for printed books works very well today, transcribing full-page handwritten mathematics exams remains unsolved because of three major roadblocks:

\subsection{The 1D Sequential Assumption of Traditional OCR}
Classic OCR engines like Tesseract \cite{smith2007overview} assume that a page consists of straight, evenly spaced lines of text. When presented with handwritten math, they try to slice tall matrix brackets into pieces, mistake fraction bars for underline noise or stray dashes, and scramble exponents across separate lines. The resulting output is usually garbled and uncompilable.

\subsection{The Over-Correction Trap of Commercial AI Models}
When Large Multimodal Models (LMMs) like GPT-4o and Claude 3.5 Sonnet became popular, it was tempting to ask them to transcribe exam pages directly. But when researchers tested this in grading settings, they ran into a dangerous flaw: \textbf{silent semantic over-correction} \cite{chen2024hallucination, pink2026overcorrection}.

Commercial models are trained to be helpful problem solvers. When they read a student's answer that contains an algebra mistake---like $(-2) \times 3 = 6$ or a wrong row operation---their language model kicks in and quietly ``fixes'' the derivation. The model outputs clean, correct \LaTeX\ math that the student never actually wrote on the page.

Recent papers investigating this behavior (such as the PINK benchmark \cite{pink2026overcorrection}) show that commercial multimodal models silently overwrite student mistakes with correct steps in $42\%$ to $66\%$ of cases. In an exam evaluation setting, this is disastrous: a student who made a mistake could receive full marks because the AI rewrote their answer. An assistive transcription tool must follow one strict rule:
\begin{equation}
\label{eq:fidelity_invariant}
\hat{\mathbf{Y}} = \arg\max_{\mathbf{Y}} P(\mathbf{Y} \mid \mathbf{I}; \Theta) \quad \text{subject to verbatim transcription fidelity.}
\end{equation}

\subsection{Student Privacy and On-Campus Supercomputing}
University answer booklets are confidential academic records containing student roll numbers, names, and grades. Uploading thousands of unencrypted exam scans to commercial third-party cloud APIs (such as OpenAI or Mathpix \cite{mathpix2020technology}) violates university data privacy policies. On top of privacy concerns, paying per-page API fees for 13,000+ pages every semester is expensive. Any practical solution must run completely on-campus, inside the university's private network. At IIT Guwahati, this means deploying our models on the \textbf{PARAM Kamrupa Supercomputer} in an isolated, air-gapped environment with no external internet connection.

\section{Project Scope: Assistive Transcription, Not Automatic Grading}
\label{sec:scope_boundaries}

We need to make our project scope crystal clear:
\begin{definition}[Project Scope: Assistive Transcription]
This project builds an \textbf{Assistive Document Transcription Pipeline}. Its only job is to convert raw scanned handwritten exam pages $\mathbf{I}$ into valid, compilable \LaTeX\ source code $\mathbf{Y}$. The system \textbf{does not grade the exam, does not evaluate mathematical correctness, and does not assign marks.}
\end{definition}

Grading remains entirely the responsibility of human professors and teaching assistants. Our software simply removes the eye-strain of reading messy handwriting by giving evaluators clean, compiled mathematics right next to the original student scan.

\subsection{Two-Semester Plan (Phase 1 vs. Phase 2)}
Our B.Tech Project is split across two semesters:
\begin{enumerate}
    \item \textbf{Semester 1 (MA4998P Project-I / Phase 1 - Current Report):} Setting up our data pipeline, organizing the 13,146 scanned pages, deploying models inside the air-gapped PARAM Kamrupa cluster, benchmarking baseline pipelines (our Modular Hybrid Pipeline and zero-shot Qwen2.5-VL-7B), and identifying key hardware bottlenecks (16GB VRAM limits on Tesla V100) and model failure modes (TrOCR prose drift, the repetition penalty paradox, and the problem-solver bias).
    \item \textbf{Semester 2 (MA4999P Project-II / Phase 2 - Future Work):} Fine-tuning Qwen2.5-VL directly on handwritten exam crops using LoRA, using INT4 quantization to process full-resolution 2000px images within VRAM limits, and building a side-by-side web review interface for instructors.
\end{enumerate}

\section{Summary of Our Contributions}
\label{sec:contributions}

The main engineering and research contributions we made during Phase 1 include:
\begin{enumerate}
    \item \textbf{Curating the MA102 Benchmark Dataset:} We organized 13,146 scanned answer pages from 939 students who took MA102 at IIT Guwahati, establishing student-disjoint splits and compiling a gold-standard benchmark transcribed verbatim by hand.
    \item \textbf{Building the Modular Hybrid Pipeline:} We designed a modular pipeline combining Sauvola adaptive binarization, projection deskewing, 1D vertical IoU line grouping, an algorithmic matrix parser (\texttt{MatrixGridParser}), a dual-head dispatcher (TrOCR for text and a custom DenseNet HMER model for formulas), and 2D DAG layout reordering.
    \item \textbf{Assistant Prefill Locking for VLMs:} We introduced an Assistant Prefill Locking technique using \texttt{continue\_final\_message=True}, which stops instruction-tuned vision-language models (Qwen2.5-VL) from generating chatty conversational summaries or Python code.
    \item \textbf{Compilation Invariant Guarantee Engine:} We formulated a formal Deterministic Pushdown Automaton (DPDA) stack balancer that automatically tracks and closes unclosed \LaTeX\ environments, guaranteeing that $100\%$ of generated documents compile into PDFs without manual fixing.
    \item \textbf{Air-Gapped Deployment on PARAM Kamrupa:} We got large neural networks running on NVIDIA Tesla V100 GPUs on PARAM Kamrupa without internet access, fixing legacy CentOS 7 C++ library mismatches through runtime \texttt{ctypes} preloading and configuring Cutlass memory-efficient attention for Volta GPUs.
    \item \textbf{Master Multi-Band High-Resolution Slicing Architecture:} We engineered a 3-band overlapping slicing pipeline with an $8\%$ vertical buffer ($\approx 180\,\text{px}$) that feeds high-resolution horizontal strips ($\text{max\_dim} = 1200\,\text{px}$) into the vision-language model, collapsing peak ViT attention memory from $1.30\,\text{GiB}$ down to $\approx 240\,\text{MiB}$ on Tesla V100 (16GB VRAM) and eliminating CUDA Out-of-Memory crashes completely.
    \item \textbf{Empirical Breakthrough and Failure Diagnosis:} We benchmarked authentic student exam proofs on PARAM Kamrupa, achieving \textbf{88.24\% Word-for-Word Token Accuracy} on complete MA102 vector subspace proofs (Roll 240101079), while diagnosing and resolving critical failure modes including the ``Over-Helpful Tutor'' prior and blank-page hallucination loops via a Laplacian edge energy filter.
\end{enumerate}

\section{Mathematical Problem Formulation}
\label{sec:task_formulation}

We treat full-page exam transcription as conditional sequence generation under strict compilation constraints:
\begin{itemize}
    \item \textbf{Input Image:} An unruled exam page is digitized as an image tensor $\mathbf{I} \in \mathbb{R}^{H \times W \times C}$ (roughly $2338 \times 1653$ pixels at 200--300 DPI).
    \item \textbf{Output Sequence:} Given a discrete \LaTeX\ vocabulary $\mathcal{V}_{\text{LaTeX}}$, the target document is an ordered sequence of tokens $\mathbf{Y} = (y_1, y_2, \dots, y_T)$. Our model computes the conditional probability:
    \begin{equation}
    P(\mathbf{Y} \mid \mathbf{I}; \Theta) = \prod_{t=1}^T P(y_t \mid y_{<t}, \mathbf{I}; \Theta)
    \end{equation}
    \item \textbf{Compilation and Verbatim Fidelity:} The predicted sequence $\hat{\mathbf{Y}}$ must compile without errors ($\text{Compile}(\hat{\mathbf{Y}}) = 1$) and transcribe the student's physical ink verbatim:
    \begin{equation}
    \hat{\mathbf{Y}} = \arg\max_{\mathbf{Y}} P(\mathbf{Y} \mid \mathbf{I}; \Theta) \quad \text{subject to } \text{Compile}(\hat{\mathbf{Y}}) = 1 \text{ and verbatim fidelity.}
    \end{equation}
\end{itemize}

\section{Evaluation Metrics}
\label{sec:evaluation_metrics}

Measuring the quality of handwritten math transcription requires more than standard character error rates. We track four quantitative metrics:
\begin{enumerate}
    \item \textbf{Exact Match Rate (EMR / $\text{ExpRate}_0$):} The percentage of expressions or blocks transcribed with zero token errors:
    \begin{equation}
    \text{EMR} = \text{ExpRate}_0 = \frac{1}{N} \sum_{i=1}^N \mathbb{I}\left( \hat{\mathbf{Y}}^{(i)} = \mathbf{Y}_{\text{ref}}^{(i)} \right)
    \end{equation}
    We also report $\text{ExpRate}_1$ and $\text{ExpRate}_2$ (allowing at most 1 or 2 token errors), following the CROHME competition standard \cite{mouchere2016icfhr}.
    \item \textbf{Token Error Rate (TER):} The normalized Levenshtein distance measured over \LaTeX\ tokens rather than raw characters:
    \begin{equation}
    \text{TER} = \frac{\sum_{i=1}^N \mathcal{D}_{\text{Lev}}\left( \hat{\mathbf{Y}}^{(i)}, \mathbf{Y}_{\text{ref}}^{(i)} \right)}{\sum_{i=1}^N |\mathbf{Y}_{\text{ref}}^{(i)}|} \times 100\%
    \end{equation}
    This prevents long commands like \texttt{\textbackslash textbf} from skewing error counts relative to single variables.
    \item \textbf{Normalized Edit Distance (NED):} A bounded score between $0$ and $1$:
    \begin{equation}
    \text{NED}\left( \hat{\mathbf{Y}}, \mathbf{Y}_{\text{ref}} \right) = 1 - \frac{\mathcal{D}_{\text{Lev}}\left( \hat{\mathbf{Y}}, \mathbf{Y}_{\text{ref}} \right)}{\max\left( |\hat{\mathbf{Y}}|, |\mathbf{Y}_{\text{ref}}| \right)}
    \end{equation}
    \item \textbf{1D Vertical IoU for Line Grouping:} The vertical overlap between two candidate line intervals $A = [y_{A1}, y_{A2}]$ and $B = [y_{B1}, y_{B2}]$:
    \begin{equation}
    \text{IoU}_{1D}(A, B) = \frac{\text{Length}(A \cap B)}{\text{Length}(A \cup B)} = \frac{\max\left(0, \ \min(y_{A2}, y_{B2}) - \max(y_{A1}, y_{B1})\right)}{\max(y_{A2}, y_{B2}) - \min(y_{A1}, y_{B1})}
    \end{equation}
    When $\text{IoU}_{1D}(A, B) \ge \theta_{\text{IoU}} = 0.35$, our line segmenter merges the two candidate strips.
\end{enumerate}

\section{Report Organization}
\label{sec:organization}

The rest of this report is organized as follows:
\begin{itemize}
    \item \textbf{Chapter \ref{ch:literature_review} (Literature Review and Foundations)} traces how mathematical expression recognition developed over time, examines attention-based sequence models, explains why our early experiment with GAN font normalization failed (Pivot 1), formalizes context-free grammars and pushdown automata for syntactic closure, and discusses layout analysis and vision-language architectures.
    \item \textbf{Chapter \ref{ch:dataset} (The MA102 Examination Dataset)} describes the 13,146 scanned exam pages from IIT Guwahati, categorizes linear algebra and multivariable calculus problems in the corpus, outlines physical ink degradations, details key student evaluation scripts, and defines our student-disjoint split.
    \item \textbf{Chapter \ref{ch:methodology} (Proposed Methodology)} gives the full technical details of our three architectural paradigms: the Modular Dual-Head Hybrid Pipeline, the Monolithic Vision-Language Engine, and our master Multi-Band High-Resolution Slicing and Pushdown Automaton Stitching Architecture.
    \item \textbf{Chapter \ref{ch:supercomputing} (Supercomputing Deployment)} details our setup on the PARAM Kamrupa Supercomputer, covering offline weight caching, our \texttt{ctypes} dynamic linking workaround for CentOS 7, Volta sm\_70 attention kernels, and VRAM memory budgeting on 16GB Tesla V100 GPUs.
    \item \textbf{Chapter \ref{ch:experiments} (Experimental Results and Conclusion)} presents quantitative benchmarks (including our 88.24\% token accuracy milestone on Student Roll 240101079), analyzes why baseline models stumble on exam scripts, summarizes our Phase 1 accomplishments, and outlines our Phase 2 plan for Semester 2.
\end{itemize}
